% !TEX program = xelatex
\documentclass[11pt]{ctexart}
%=====================================================================
% theme-violet.tex —— 讲义精讲 · 紫花主题
% 用法：\documentclass[11pt]{ctexart}
%       \input{assets/theme-violet}
%       \begin{document} ... \end{document}
% 只用 ctexart 自带能力 + xcolor/enumitem/fancyhdr/pifont/amsmath，
% 不依赖 tikz/tcolorbox，保证 MiKTeX 可直接编译。
%=====================================================================

\usepackage[a4paper,left=2.0cm,right=2.0cm,top=2.1cm,bottom=2.3cm,%
            headsep=12pt,footskip=22pt]{geometry}
\usepackage{amsmath,amssymb}
\usepackage{needspace}
\usepackage{xcolor}
\usepackage{array}
\usepackage{enumitem}
\usepackage{pifont}
\usepackage{fancyhdr}
\usepackage[hidelinks]{hyperref}

%---------------------------------------------------------------- 配色
\definecolor{violet}{RGB}{107,78,143}      % 主色：紫
\definecolor{violetdeep}{RGB}{63,44,92}    % 深紫：标题/强调
\definecolor{petal}{RGB}{183,159,208}      % 花瓣：装饰线、花饰
\definecolor{bandviolet}{RGB}{247,243,251} % 浅紫底：档案/提示框
\definecolor{bandcream}{RGB}{252,249,242}  % 米白底：易错/注意框
\definecolor{rulepurple}{RGB}{217,207,230} % 分隔线
\definecolor{inkgray}{RGB}{92,96,110}      % 次要文字

%---------------------------------------------------------------- 字体
% 有 Noto CJK 用 Noto，没有就退回 ctex 默认（Windows 下为宋体/黑体）
\IfFontExistsTF{Noto Serif CJK SC}{\setCJKmainfont{Noto Serif CJK SC}}{}
\IfFontExistsTF{Noto Sans CJK SC}{\setCJKsansfont{Noto Sans CJK SC}}{}

%---------------------------------------------------------------- 版式
\setlist{nosep,leftmargin=1.9em}
\linespread{1.28}
\setlength{\parskip}{3.5pt}
\setlength{\parindent}{0pt}
\setlength{\tabcolsep}{5pt}
\renewcommand{\arraystretch}{1.28}

%---------------------------------------------------------------- 花饰
% \ding{95} = 五瓣小花，\ding{96} = 星形花；都来自 pifont，必定存在
\newcommand{\orn}{\textcolor{petal}{\ding{95}}}          % 行内小标
\newcommand{\ornstar}{\textcolor{petal}{\ding{96}}}      % 星形花
\newcommand{\flowerrule}{%                                % 居中花饰分隔线
  \par\vspace{7pt}%
  \noindent\makebox[\linewidth]{\color{petal}%
    \rule{0.20\linewidth}{0.6pt}\hspace{8pt}{\small\ding{95}}\hspace{4pt}%
    \raisebox{1.1pt}{\tiny\ding{95}}\hspace{4pt}{\small\ding{95}}\hspace{8pt}%
    \rule{0.20\linewidth}{0.6pt}}%
  \par\vspace{7pt}}

%---------------------------------------------------------------- 标题与节头
\newcommand{\doctitle}[3]{%   #1 主标题  #2 副标题  #3 元信息（可用 \\ 分行）
  \begin{center}
    {\color{petal}\small \orn\enspace #2\enspace \ornstar}\\[7pt]
    {\LARGE\bfseries\color{violetdeep} #1}\\[10pt]
    {\color{petal}\rule{0.55\linewidth}{1pt}}\\[4pt]
    {\small\color{inkgray} #3}
  \end{center}\vspace{2pt}}

\newcommand{\hdr}[2]{%   #1 章节号（如 §0）  #2 标题
  \par\vspace{12pt}%
  \noindent\colorbox{violetdeep}{\color{white}\bfseries\small\;#1\;}\hspace{7pt}%
  {\color{violetdeep}\large\bfseries #2}\par\nobreak
  \vspace{3pt}{\color{petal}\hrule height 1.2pt}\vspace{6pt}}

\newcommand{\subhdr}[1]{%
  \par\vspace{6pt}%
  \noindent{\color{violet}\bfseries #1}\par\vspace{1pt}}

%---------------------------------------------------------------- 提示框
\newcommand{\callout}[1]{% 浅紫底：讲义档案、读法说明
  \par\medskip\noindent
  \fcolorbox{rulepurple}{bandviolet}{\parbox{\dimexpr\linewidth-2\fboxsep-2\fboxrule}{\vspace{2pt}#1\vspace{2pt}}}%
  \par\medskip}

\newcommand{\warnbox}[1]{% 米白底：易错、边界条件
  \par\medskip\noindent
  \fcolorbox{rulepurple}{bandcream}{\parbox{\dimexpr\linewidth-2\fboxsep-2\fboxrule}{\vspace{2pt}#1\vspace{2pt}}}%
  \par\medskip}

\newcommand{\gapitem}[2]{% #1 缺口类型（紫） #2 内容
  \par\vspace{2pt}\noindent
  \hangindent=1.4em\hangafter=1
  {\color{petal}\small\ding{96}}\enspace{\bfseries\color{violetdeep}#1}\enspace #2\par}

% 题干块（\qstem）：每道例题／作业的「详解」开头必须先复述题干与出处。
% 起因：2026-10-25 第五章讲义把七道作业的题面集中放在一页、详解另起一页，
% 读者翻到「第 4 题」时看不到条件，等于一道没有题干的题目。
% 规则：详解块内自带题干，读者不翻页就能读懂题在问什么。
\newcommand{\qstem}[1]{%
  \par\medskip\noindent
  \fcolorbox{violet}{bandviolet}{%
    \parbox{\dimexpr\linewidth-2\fboxsep-2\fboxrule}{%
      \vspace{2pt}%
      {\small\bfseries\color{violet}题干}\hspace{0.8em}#1%
      \vspace{2pt}}}%
  \par\medskip}

% v2: 保护题干与开头，长题可跨页，续页补上下文。
\newcommand{\qbreak}{\needspace{6\baselineskip}}
% 独立练习：ID、来源、题干、考点、答案标签（通常 ans:ID）。
\newcommand{\practice}[5]{%
  \qbreak\subhdr{练习 #1}\qstem{（#2）#3}%
  \textbf{考点}：#4\par
  详解见附录第~\pageref{#5}~页。\par\medskip}
% 附录开头：ID、来源、条件摘要；自动产生 ans:ID 标签。
\newcommand{\answerstart}[3]{%
  \qbreak\subhdr{#1\quad 详解}\label{ans:#1}\qstem{（#2）#3}}
\newcommand{\answercontinue}[2]{%
  \needspace{5\baselineskip}\subhdr{#1\quad 续解}%
  {\small\color{inkgray}本页接续：#2}\par\medskip}
% 延迟提供默认值，让旧讲义可在导言区保留自己的 \newcommand{\step}。
\AtBeginDocument{\providecommand{\step}[1]{\needspace{3\baselineskip}\par\medskip\textbf{#1}\par}}

% 前提旁注：只放在「需要理解的那一步」后面，给陈述 + 本题代入，不重证。
\newcommand{\aside}[1]{%
  \par\vspace{1pt}\noindent
  \hangindent=1.6em\hangafter=1
  {\small\color{inkgray}%
    \hspace*{0.6em}{\color{petal}$\triangleright$}\hspace{2pt}#1}%
  \par\vspace{2pt}}

%---------------------------------------------------------------- 表格
% 本机 MiKTeX 的 colortbl 与 array 版本不匹配（\rowcolor 会报错），
% 故表格统一用「上粗紫线 + 表头 + 中细线」的三线式，不依赖 colortbl。
\newcommand{\tabtop}{\par\vspace{3pt}{\color{violet}\hrule height 1.2pt}\vspace{1pt}}
\newcommand{\tabmid}{\par\vspace{1pt}{\color{petal}\hrule height 0.5pt}\vspace{1pt}}
\newcommand{\tabbot}{\par\vspace{3pt}{\color{rulepurple}\hrule height 0.5pt}\vspace{8pt}}
\newcommand{\thead}[1]{{\bfseries\color{violetdeep}#1}}

%---------------------------------------------------------------- 页眉页脚
\pagestyle{fancy}
\fancyhf{}
\renewcommand{\headrulewidth}{0.4pt}
\renewcommand{\headrule}{\hbox to\headwidth{\color{rulepurple}\leaders\hrule height \headrulewidth\hfill}}
\fancyhead[L]{\small\color{inkgray}\songti \doctag}
\fancyhead[R]{\small\color{inkgray}\songti \coursetag}
\fancyfoot[C]{\small\color{inkgray}\orn\enspace{\color{violet}\thepage}\enspace\orn}

% 每篇讲义在导言区设这两个值
\newcommand{\doctag}{}
\newcommand{\coursetag}{}
\newcommand{\settags}[2]{\renewcommand{\doctag}{#1}\renewcommand{\coursetag}{#2}}

%---------------------------------------------------------------- 数学记号
\newcommand{\vect}[1]{\boldsymbol{#1}}
\newcommand{\dd}{\,\mathrm{d}}
\newcommand{\ddx}{\frac{\mathrm{d}}{\mathrm{d}x}}
\addtolength{\jot}{2pt}



\usepackage{graphicx}
\setlength{\headheight}{22pt}
\settags{实验 5 · 电路图页}{电工电子技术}

\begin{document}

\begin{center}
{\LARGE\bfseries 实验 5\quad 共射极单管放大电路\quad 电路图页}\\[4pt]
{\small 图 5.2.4（教材 p.151 原图）· 端子与开关对照 · 报告里怎么用\ \textperiodcentered\ 2026-11-05}
\end{center}

\begin{center}
\includegraphics[width=0.66\linewidth]{图-实验5-教材图524-无题注.png}

{\small 图 5.2.4\quad 共射极单管放大电路（\textbf{裁剪}自教材书 p.151）。图中 a、b1、b2、d1、d2、c 六处小圆圈是\textbf{实验箱面上的接线柱}，不是元件；$S_1$、$S_2$ 是箱面开关，$R_{P1}$、$R_{P2}$ 是箱面电位器。}
\end{center}

\begin{center}
\footnotesize
\begin{tabular}{@{}p{1.7cm}p{6.2cm}p{6.6cm}@{}}
\hline
\thead{端子/开关} & \thead{它在图上的位置与作用} & \thead{箱面上怎么接、接错会怎样} \\
\hline
a & $R_{P1}$ 的下端，与 b1 或 b2 之间跳线 & 跳线选上偏置电阻；悬空则基极无偏置，$I_C\approx0$ \\
b1 & $R_{B11}=30\,\mathrm{k}\Omega$ 的上端 & 与 a 相连时用 30 k 作上偏置 \\
b2 & $R_{B12}=150\,\mathrm{k}\Omega$ 的上端 & 与 a 相连时用 150 k 作上偏置 \\
d1 & $R_{B21}=39\,\mathrm{k}\Omega$ 的下端（接地侧） & 悬空则基极电位被抬到接近 $U_{CC}$，管子饱和 \\
d2 & $R_{B22}=12\,\mathrm{k}\Omega$ 的下端（接地侧） & 同上 \\
c & 公共地（与电源负极、示波器地夹同一根线） & 多点乱接地会引入干扰，读数抖动 \\
$S_1$ & 与 $R_{E2}=200\,\Omega$ 并联 & \textbf{闭合＝短接 $R_{E2}$＝没有这级负反馈}；记反了表 5.2.3 全反 \\
$S_2$ & 串在 $R_{L1}=3\,\mathrm{k}\Omega$ 负载支路里 & 闭合＝带载，断开＝开路；开路时测的是 $U_{oc}$ \\
$R_{P1}$ & $10\,\mathrm{k}\Omega$，静态微调 & 静态测量时调到零（$0\,\Omega$），四组偏置才可比 \\
$R_{P2}$ & $10\,\mathrm{k}\Omega$，负载细调 & 把 $R_L$ 从 $3\,\mathrm{k}\Omega$ 调到 $13\,\mathrm{k}\Omega$，调完要量实际阻值 \\
\hline
\end{tabular}
\end{center}

\warnbox{\textbf{通电前必须用万用表核对的四件事}：（一）a 与 b1/b2 的跳线是否按偏置组合接对；（二）基极到地的阻值是否约为 $R_{B1}+R_{P1}$；（三）$C_E$（$100\,\mu\mathrm{F}$）两端是否只跨 $R_{E1}=1\,\mathrm{k}\Omega$（若发现它跨在 $R_{E1}+R_{E2}$ 上，静态点与 $A_u$ 都会对不上算例）；（四）$C_1$、$C_2$、$C_E$ 三只电解电容的 $+$ 端是否朝高电位（$C_1$ 朝基极、$C_2$ 朝集电极、$C_E$ 朝发射极）。}

\subhdr{报告里怎么用这张图（三条，照做即可）}

\begin{enumerate}
  \item \textbf{引用句（英文，可直接抄进 II. Experimental Principles）}：\textit{The common-emitter amplifier used in this experiment is shown in Fig.~1 (textbook Fig.~5.2.4). The base bias is supplied by the divider $R_{B1}$/$R_{B2}$; $R_{P1}$ was set to zero during the static measurements; $S_1$ shorts $R_{E2}$ when closed, and $S_2$ connects the load $R_{L1}$.}
  \item \textbf{手绘顺序（五分钟能画完）}：先画 $+U_{CC}$ 与地两条横线 → 画晶体管 T（NPN，箭头朝外）→ 集电极串 $R_C=3\,\mathrm{k}\Omega$ → 基极接分压 $R_{B1}$、$R_{B2}$，上端经 $R_{P1}$ 到 $+U_{CC}$ → 发射极串 $R_{E1}$、$R_{E2}$，$C_E$ 并在 $R_{E1}$ 两端、$S_1$ 并在 $R_{E2}$ 两端 → 输入端 $C_1$ 串 $u_i$，输出端 $C_2$ 串负载支路（$S_2$、$R_{L1}$、$R_{P2}$）。
  \item \textbf{图上必须标出的量}：$R_{B11}=30\,\mathrm{k}$、$R_{B12}=150\,\mathrm{k}$、$R_{B21}=39\,\mathrm{k}$、$R_{B22}=12\,\mathrm{k}$、$R_C=3\,\mathrm{k}$、$R_{E1}=1\,\mathrm{k}$、$R_{E2}=200\,\Omega$、$R_{P1}=R_{P2}=10\,\mathrm{k}$、$C_1=C_2=27\,\mu\mathrm{F}$、$C_E=100\,\mu\mathrm{F}$、$\beta=100$、$U_{CC}=12\,\mathrm{V}$。少标一个，批阅时会被要求补。
\end{enumerate}

\callout{\textbf{一句话记住这张图}：$R_{P1}$ 选上偏置、b1/b2 选 30 k 或 150 k、d1/d2 选 39 k 或 12 k——四组偏置就是这四只电阻的两种两种组合；$S_1$ 管负反馈、$S_2$ 管负载，其余全是 $\pm12\,\mathrm{V}$、地、输入输出四根线。}

\end{document}
